\begin{code}[hide]
module vqupit.sections.warmup where

open import Data.Nat using (ℕ ; suc)
open import Data.Product using (_×_ ; _,_ ; proj₁ ; proj₂)
open import Data.Sum using (inj₁ ; inj₂)
open import Data.Unit using (⊤ ; tt)
open import Notations using (₁₊ ; ₂₊ ; ₃₊)
open import Word.Base using (Word ; WRel ; [_]ʷ ; ε ; _•_ ; wmap)
open import Presentation.Definitions using (_IsPresentationOf_)
open import Presentation.Construct.Base using (_⊎^_ ; _⊕^_ ; _⋄_⋄_ ; ConjRelʷ)
open import Examples.Groups.Cyclic.Normalization using (_Cn,_===_)

postulate
  proof-elided : ∀ {ℓ} {A : Set ℓ} → A

private variable
  n : ℕ

data Gate : ℕ → Set where
  σ-gate : Gate 2
data Gen : ℕ → Set where
  gate₁ : Gate 1 → Gen (₁₊ n)
  gate₂ : Gate 2 → Gen (₂₊ n)
  _↥   : Gen n → Gen (₁₊ n)

Circuit : ℕ → Set
Circuit n = Word (Gen n)

_↑ : Circuit n → Circuit (₁₊ n)
_↑ = wmap _↥

pattern σ-gen = gate₂ σ-gate

σ : Circuit (₂₊ n)
σ = [ σ-gen ]ʷ

infix 4 _SRel,_===_ _VRel,_===_ _≈_
infixr 10 σ•_

postulate
  _≈_ : Circuit n → Circuit n → Set
  _VRel,_===_ : (n : ℕ) → WRel (Gen n)
  sw₂ : ∀ {n} → ⊤ ⊎^ n → ⊤ ⊎^ n

module Wreath where
  postulate wreath-group : ℕ → ℕ → Set
\end{code}
\section{Warm-up: Permutations and Signed Permutations}\label{sec:warmup}

Two small examples from the library exercise everything the symplectic
development needs.  Swap circuits, which present the symmetric groups,
show a coset tower whose obligations close by evaluation; the wreath
products $\mathbb{Z}_N \wr S_n$, the groups of signed permutations, show
the semidirect-product theorem gluing two presentations along a
syntactic action.  The symplectic factor of \S\ref{sec:symplectic} is
the first example with a richer action, and the Clifford group of
\S\ref{sec:composing} is the second example with $\Sp$ in place of
$S_n$.

\subsection{Swap circuits and the symmetric groups}\label{sec:sn}

Permutation circuits have one gate, the swap $\sigma$ on two wires, and
two group-specific axioms, stated once each with the gates at the bottom
and the width left open, so that \aid{cong↑} moves them to any height:

\begin{minipage}[c]{0.68\textwidth}
\begingroup\renewcommand{\AgdaCodeStyle}{\footnotesize}
\begin{code}
data _SRel,_===_ : (n : ℕ) → WRel (Gen n) where
  order       : (₂₊ n) SRel, σ • σ === ε
  yang-baxter : (₃₊ n) SRel, σ • σ ↑ • σ === σ ↑ • σ • σ ↑
\end{code}
\endgroup
\end{minipage}\hfill
\begin{minipage}[c]{0.28\textwidth}\centering
  \ufig[0.45]{sn-order}\\[3pt]
  \ufig[0.45]{sn-yb}
\end{minipage}

\medskip\noindent These are two thirds of the Coxeter presentation of
$S_n$; the third Coxeter relation, that swaps on disjoint wires commute,
is the structural rule \aid{comm₂} and comes for free.  The coset chain
is $S_1 < S_2 < \dots < S_{n+1}$, where $S_n$ sits inside $S_{n+1}$ as
the circuits of the form $c{\uparrow}$.  Two circuits lie in the same
right coset of $S_n$ exactly when they bring the same wire to the
bottom, so $S_n$ has index $n{+}1$, and a natural representative of each
coset is the \emph{descending staircase} of swaps that carries that
wire down.  A coset label is the length of the staircase, and the
section maps it to the staircase circuit:

\begin{minipage}[c]{0.36\textwidth}
\begingroup\renewcommand{\AgdaCodeStyle}{\footnotesize}
\begin{code}
data C : ℕ → Set where
  ε   : C n
  σ•_ : C n → C (₁₊ n)

[_]ᶜ : C n → Circuit (₁₊ n)
[ ε ]ᶜ     = ε
[ σ• c ]ᶜ  = σ • ([ c ]ᶜ ↑)
\end{code}
\endgroup
\end{minipage}\hfill
\begin{minipage}[c]{0.6\textwidth}\centering\small
\begin{tabular}{@{}l@{\;\;}ccc@{}}
  $\aid{C}\;2$: & \ufig[0.45]{c2-0} & \ufig[0.45]{c2-1} & \ufig[0.45]{c2-2}\\
  & $\varepsilon$ & $\sigma{\bullet}\:\varepsilon$ &
    $\sigma{\bullet}\:\sigma{\bullet}\:\varepsilon$\\[3pt]
  $\aid{C}\;1$, $\aid{C}\;0$: & \ufig[0.45]{c1-0} & \ufig[0.45]{c1-1} & \ufig[0.45]{c0-0}\\
  & $\varepsilon$ & $\sigma{\bullet}\:\varepsilon$ & $\varepsilon$
\end{tabular}
\end{minipage}

\medskip\noindent Recursing on the coset writes every $f \in S_{n+1}$
uniquely as $c_1 c_2 \cdots c_n$ with $c_k \in \aid{C}\;k$, a numeral
in the factorial number system.  \emph{Computing} this factorisation is
the coset table: given a staircase and one generator, it says where
their product lands (the updated coset) and what is left over on the
wires above (the residual circuit), so that staircase followed by
generator equals residual, shifted up, followed by updated staircase.
Pushing $\sigma$ into a staircase can lengthen it, cancel with it, or
pass through by Yang--Baxter; a generator that does not touch the
bottom wire commutes past or recurses (Figure~\ref{fig:pushing}):

\begin{minipage}[c]{0.7\textwidth}
\begingroup\renewcommand{\AgdaCodeStyle}{\scriptsize}
\begin{code}
ract : C (₁₊ n) → Gen (₂₊ n) → Circuit (₁₊ n) × C (₁₊ n)
ract ε         σ-gen = ε , σ• ε                       -- grows
ract (σ• ε)    σ-gen = ε , ε                          -- cancels
ract (σ• σ• c) σ-gen = σ , σ• σ• c                    -- passes
ract ε         (g ↥) = [ g ]ʷ , ε                     -- commutes past
ract (σ• c)    (g ↥) = proj₁ (ract c g) ↑ , σ• (proj₂ (ract c g))

ract-sound : ∀ c b → let (b' , c') = ract c b
                     in  [ c ]ᶜ • [ b ]ʷ ≈ b' ↑ • [ c' ]ᶜ
\end{code}
\endgroup
\begin{code}[hide]
ract-sound = proof-elided
\end{code}
\end{minipage}\hfill
\begin{minipage}[c]{0.27\textwidth}\centering
  \ufig[0.45]{ract-l}\;\scalebox{0.7}{$\approx$}\;\ufig[0.45]{ract-r}
\end{minipage}

\begin{figure}
\centering
\begin{minipage}[c]{0.72\textwidth}\footnotesize
\begin{tabular}{@{}r@{\,}c@{\,}l@{\;}l@{\quad}r@{\,}c@{\,}l@{\;}l@{}}
  \ufig[0.4]{push4-l} & \scalebox{0.7}{$\rightarrow$} & \ufig[0.4]{push4-r} & grows &
  \ufig[0.4]{push3-l} & \scalebox{0.7}{$\rightarrow$} & \ufig[0.4]{push3-r} & cancels\\[3pt]
  \ufig[0.4]{push2-l} & \scalebox{0.7}{$\rightarrow$} & \ufig[0.4]{push2-r} & passes &
  \ufig[0.4]{push1-l} & \scalebox{0.7}{$\rightarrow$} & \ufig[0.4]{push1-r} & commutes past
\end{tabular}
\end{minipage}\hfill
\begin{minipage}[c]{0.24\textwidth}\centering
  \ufig[0.55]{sn-stair}
\end{minipage}
\caption{Left: pushing a swap through a staircase, the four cases of
the coset table \aid{ract}, in the order of its clauses.  Right: the
section of a normal form in $\aid{NF}\;4$, one staircase per level.}
\label{fig:pushing}
\end{figure}

\noindent The residual is empty or a single swap here, but for richer
gate sets it can be a longer circuit, which is why the table returns a
circuit one wire up rather than a generator.  The soundness law
\aid{ract-sound} is hypothesis (5) of \S\ref{sec:engine}, proved by a
short derivation per clause; the other four hypotheses close by
evaluation.  The tower's carrier is $\aid{NF}\;(k{+}1) = \aid{NF}\;k
\times \aid{C}\;k$, so $\aid{NF}\;n$ has $n!$ elements, and the section
concatenates the staircases, suitably shifted (Figure~\ref{fig:pushing},
right).  Uniqueness against the permutation semantics is an induction on
the tower, since equal denotations force equal top digits; soundness
and uniqueness feed \aid{by-normalization}, and the two axioms plus the
structural rules present the permutation group of $\mathsf{Fin}\;n$.

\subsection{The wreath product \texorpdfstring{$\mathbb{Z}_N \wr S_n$}{Z_N wr S_n} as a semidirect product}\label{sec:snd}

The second example composes two presentations instead of building a
tower.  Fix a base order $N = m{+}1$ and $n$ wires.  The cyclic
presentation $\langle T \mid T^{N} = \varepsilon\rangle$, iterated by the
$n$-fold direct-product construction $\Gamma_0 \aid{⊕\char94} n$,
presents $(\mathbb{Z}_N)^n$ over the alphabet $\top \aid{⊎\char94} n$
of $n$ copies of $T$, one per coordinate; the swap circuits of
\S\ref{sec:sn} present $S_n$.  The wreath product $\mathbb{Z}_N \wr S_n
= (\mathbb{Z}_N)^n \rtimes S_n$ is $S_n$ acting on the coordinates, and
the action is specified \emph{syntactically}, generator by generator:
the bottom swap exchanges the bottom two coordinates, and a shifted swap
acts on the shifted coordinates and fixes the bottom one.

\begin{code}
act : ∀ {n} → Gen n → ⊤ ⊎^ n → ⊤ ⊎^ n
act        (gate₂ σ-gate) x        = sw₂ x
act {₂₊ n} (g ↥)          (inj₁ tt) = inj₁ tt
act {₂₊ n} (g ↥)          (inj₂ y)  = inj₂ (act g y)

conj : ∀ {n} → Gen n → ⊤ ⊎^ n → Word (⊤ ⊎^ n)
conj g x = [ act g x ]ʷ
\end{code}

\noindent The semidirect theorem of \S\ref{sec:constructions} asks for
two compatibility lemmas.  \aid{conj-hypn}, that the action respects
the product relations in its second argument, is a finite case split:
each generator carries $T^{N} = \varepsilon$ on one coordinate to the
same relation on another, and the commutation of two coordinates to the
commutation of their images.  \aid{conj-hyph}, that the action respects
the Coxeter relations in its first argument, reduces to the fact that
the permutations induced by the two sides of each axiom are equal as
functions on coordinates --- $\sigma^2$ is the identity and both sides
of Yang--Baxter are the same 3-cycle --- after which the two conjugates
are propositionally equal words.  Feeding the two factor presentations
and the two lemmas into the theorem gives

\begingroup\renewcommand{\AgdaCodeStyle}{\footnotesize}
\begin{code}
wreath-presentation : ∀ n m → (((suc m Cn,_===_) ⊕^ n) ⋄ (n VRel,_===_) ⋄ ConjRelʷ conj)
                                IsPresentationOf (Wreath.wreath-group n m)
\end{code}
\endgroup
\begin{code}[hide]
wreath-presentation = proof-elided
\end{code}

\noindent where $\aid{wreath-group}\;n\;m$ is the semantic
$(\mathbb{Z}_N)^n \rtimes S_n$ built by the library's semidirect
product from the cyclic groups and the permutation group.  The lifted
normal form is a pair of $n$ residues modulo $N$ and the factorial tuple
of a permutation, so its section is a word of $T$-powers, one per
coordinate, followed by the staircases of \S\ref{sec:sn}.  At $N = 2$
this is the hyperoctahedral group of signed permutations; at general
$N$ it is the group $G(N,1,n)$ of monomial matrices with $N$-th roots of
unity as entries, the ``permutation-and-phase'' part of every Clifford
group.  The theorem delivers all of these presentations, for every $N$
and $n$, from two small lemmas: nothing about the wreath product was
proved by hand.  The Clifford group of \S\ref{sec:composing} is this
example with $\Sp$ in place of $S_n$ and the projective Pauli group in
place of $(\mathbb{Z}_N)^n$; there the action is richer and the two
lemmas are proved by semantics rather than by case split, but the shape
of the argument is exactly this one.
